% STATUS: written 2026-11-05 (Bresar--Chebotar--Martindale 2007 read at
% its preface and at the pages that carry the definition of a functional
% identity and the first example; excerpts in excerpts/. The book body is
% scanned, so only the OCR-transcribed pages are reported.)
\section{The theory of functional identities}\label{sec:functional-identities:fi-theory}

A functional identity is an equation with functions appearing as unknowns
\cite[p.~9]{BresarMartindaleChebotar2007FI}. One refers to functions
satisfying a certain functional identity as solutions of that identity
\cite[p.~9]{BresarMartindaleChebotar2007FI}. One calls the identically
zero functions the standard solution, and the aim of the theory is to
describe exactly when nonstandard solutions exist
\cite[p.~9]{BresarMartindaleChebotar2007FI}. The first functional identity
is the two-term identity $E(x)y+F(y)x=0$ on a ring $A$
\cite[Example~1.1, p.~10]{BresarMartindaleChebotar2007FI}. Its solvability
is a structural statement about $A$
\cite[Example~1.1, p.~10]{BresarMartindaleChebotar2007FI}. A semiprime
ring $A$ has a nonstandard solution exactly when it contains a nonzero
central ideal \cite[Example~1.1, p.~10]{BresarMartindaleChebotar2007FI}.
A prime ring has one exactly when it is commutative
\cite[Example~1.1, p.~10]{BresarMartindaleChebotar2007FI}. The matrix
ring $M_n(C)$ over a commutative unital ring has one exactly when $n=1$
\cite[Example~1.1, p.~10]{BresarMartindaleChebotar2007FI}.

The central concept upon which the general theory rests is the
\emph{$d$-free set} \cite[p.~x]{BresarMartindaleChebotar2007FI}. A ring is
$d$-free when two designated functional identities of degree $d$ have only
the standard solution \cite[p.~21]{BresarMartindaleChebotar2007FI}. The
concept carries the important weight of the whole development
\cite[p.~21]{BresarMartindaleChebotar2007FI}. The first
parts of the book introduce $d$-free sets and study their basic properties
and constructions \cite[p.~x]{BresarMartindaleChebotar2007FI}. They then
prove the existence of important classes of $d$-free sets inside prime and
semiprime rings \cite[p.~x]{BresarMartindaleChebotar2007FI}.

Until this monograph the general functional identity theory existed only
in research journals; the book was written because the theory had reached
maturity and deserved a comprehensive account
\cite[p.~ix]{BresarMartindaleChebotar2007FI}. The theory grew directly
out of the commuting and centralizing maps of
Section~\ref{sec:functional-identities:commuting-maps}: Bre\v{s}ar initiated
the study of functional identities based on his characterization of
centralizing maps, as reported at \cite[p.~333]{LeeTK1997Centralizing}.
The next section reports the central role of $d$-free sets
\cite[p.~x]{BresarMartindaleChebotar2007FI}. The two following sections
report the two families of applications, to derivation problems and to
centralizers and multipliers
\cite[p.~x]{BresarMartindaleChebotar2007FI}.
